\textbf{Background:} Digital patient frameworks are expanding rapidly, but terms such as digital twin, virtual patient, in silico patient, and synthetic patient are used inconsistently.
\textbf{Objective:} To synthesize how recent healthcare reviews define, construct, validate, and operationalize the digital patient spectrum.
\textbf{Methods:} Following the PROSPERO protocol, we searched health, engineering, and grey-literature sources for English-language reviews published from 2020 onward and extracted review-level data with a CADIMA-aligned form. A version-controlled pipeline applies documented record-level classification rules and regenerates all manuscript tables and figures.
\textbf{Results:} The reproducible analytic set contains 64 coded records representing 62 unique normalized review titles. Applications span multiple clinical domains, while validation, reproducibility, interoperability, governance, and clinical maturity remain unevenly reported.

\textbf{Conclusions:} The field is conceptually active but infrastructurally immature. Translation will require clearer taxonomy, stronger validation standards, interoperable data architecture, and governance models for clinically deployable patient digital twins.
